矩阵乘法中，一个输出分块的结果通常来自多个 K 维切片。程序先准备累加器，再逐次装入对应的 A、B 分块，把每次乘加结果留在同一个累加器中，最后写回输出。MMA 把这个过程分为四步：初始化累加器、装入操作数、计算并累加、写回结果。

\begin{itemize}[leftmargin=1.6em]
\item \co@mma.fill 0.0@：初始化累加器。
\item \co@mma.load@：装入当前分块的操作数。
\item \co@mma.row.row mc, ma, mb@：执行 C += A × B，将结果累加到 \co@mc@。
\item \co@mma.store mc, dst@：写回结果。
\end{itemize}

装入与计算步骤沿 K 维重复，每次把一个切片的结果累加到 \co@mc@；完成后再写回输出。底层的线程数据分配由编译器结合布局约定处理。

\co@row.row@ 表示两个操作数采用行优先布局。其他组合包括 \co@row.col@、\co@col.row@ 和 \co@col.col@。布局必须与输入数据一致，否则计算结果可能错误。

若 K 被分为两个切片，先初始化一次 \co@mc@，装入第一对操作数并累加，再装入第二对并累加到同一个 \co@mc@，最后才写回。装入句柄描述操作数片段，累加器则跨迭代保留部分结果，因此填充与写回应放在反复执行的装入和乘加过程之外。
